\section*{Chapter \ref{ch:g5:natural-hazards} --- Natural Hazards}

\begin{solution}{exo:g5:natural-hazards:1}
(a) A \omterm{def:g5:natural-hazards:hazard}{natural hazard} only: no one is harmed. (b) A \omterm{def:g5:natural-hazards:disaster}{natural disaster}: homes
are damaged. (c) A \omterm{def:g5:natural-hazards:hazard}{natural hazard} only. (d) A \omterm{def:g5:natural-hazards:disaster}{natural disaster}.
\end{solution}

\begin{solution}{exo:g5:natural-hazards:2}
From the air: a heat wave, a \omterm{def:g5:natural-hazards:storm}{storm}, a \omterm{def:g5:natural-hazards:drought}{drought}. From water: a \omterm{def:g5:water-shapes-the-land:flood}{flood}. From
the ground: a \omterm{def:g5:natural-hazards:landslide}{landslide}, a volcanic \omterm{def:g4:volcanoes-earthquakes-first-look:eruption}{eruption}, an \omterm{def:g4:volcanoes-earthquakes-first-look:earthquake}{earthquake}.
\end{solution}

\begin{solution}{exo:g5:natural-hazards:3}
A \omterm{def:g5:natural-hazards:drought}{drought} is a long period, of months or years, with much less \omterm{def:g4:clouds-rain-snow:precipitation}{rain} than
usual. It can leave farmers without a harvest, because the crops wilt and
die, and it can leave towns short of drinking water, because \omterm{def:g1:landscapes-around-us:water}{rivers} and
\omterm{def:g1:landscapes-around-us:water}{lakes} shrink.
\end{solution}

\begin{solution}{exo:g5:natural-hazards:4}
For example: bring in or tie down anything the wind could carry away
(garden chairs, toys, bins); check the torches and keep spare batteries and
water ready; stay indoors, away from the windows, during the \omterm{def:g5:natural-hazards:storm}{storm}; keep
listening to the radio for news.
\end{solution}

\begin{solution}{exo:g5:natural-hazards:5}
(a) A hurricane. (b) A typhoon. (c) A tropical cyclone.
\end{solution}

\begin{solution}{exo:g5:natural-hazards:6}
Green, yellow, orange, red. At the red level, there is great danger: act at
once and follow the instructions of the people in charge.
\end{solution}

\begin{solution}{exo:g5:natural-hazards:7}
Part of the hillside on the right broke away and slid down, leaving a bare
scar. Its \omterm{def:g3:stones-pebbles-sand:sand}{mud} spread right across the \omterm{def:g1:landscapes-around-us:landform}{valley} floor, covering the road. It
also blocked the \omterm{def:g1:landscapes-around-us:water}{river} like a dam: the \omterm{def:g1:landscapes-around-us:water}{river} water, which can no longer
flow away, has piled up behind it and flooded the woods upstream.
\end{solution}

\begin{solution}{exo:g5:natural-hazards:8}
Lena is right. We cannot stop the \omterm{def:g4:clouds-rain-snow:precipitation}{rain} or the \omterm{def:g1:landscapes-around-us:water}{river}, but we can keep people
out of harm's way. For example: build houses on higher ground instead of on
the \omterm{def:g5:water-shapes-the-land:flood}{floodplain}; warn people in time with \omterm{def:g5:water-shapes-the-land:flood}{flood} warnings, so that they move
to higher ground; teach everyone never to walk or drive through \omterm{def:g5:water-shapes-the-land:flood}{flood}
water.
\end{solution}

\begin{solution}{exo:g5:natural-hazards:9}
(a) $4 \times 60 = 240$ millimetres. (b) A quarter of 60 is 15; in four
months, $4 \times 15 = 60$ millimetres. (c) $240 - 60 = 180$ millimetres
were missing: three quarters of the usual \omterm{def:g4:clouds-rain-snow:precipitation}{rain}.
\end{solution}

\begin{solution}{exo:g5:natural-hazards:10}
(a) $1\,200 \div 4 = 300$ homes. (b) A third of 300 is 100; so
$300 - 100 = 200$ homes are still on the \omterm{def:g5:water-shapes-the-land:flood}{floodplain}.
\end{solution}

\begin{solution}{exo:g5:natural-hazards:11}
Possible reasons: the first village had been warned in time and people had
prepared (shutters closed, objects brought in), the second had not; the
houses of the first are solidly built, those of the second are fragile; the
second village lies on a low \omterm{def:g1:landscapes-around-us:water}{coast} or a \omterm{def:g5:water-shapes-the-land:flood}{floodplain}, where the \omterm{def:g5:natural-hazards:storm}{storm} also
brought a \omterm{def:g5:water-shapes-the-land:flood}{flood}, while the first stands on higher ground; trees close to the
houses fell on them in the second village. The same hazard does not cause
the same harm everywhere: it depends on where and how people live, and on
how well they are prepared.
\end{solution}
